\section{Introduction and Theory}

\subsection{Inverting Amplifier}

The inverting amplifier, drawn in Figure~\ref{fig:invamp}, produces an
output $V_{\text{out}}$ that is an amplified, sign-inverted copy of the
input $V_{\text{in}}$. Under the ideal op-amp assumptions the closed-loop
gain is
\begin{equation}
  A_{v} \;=\; \frac{V_{\text{out}}}{V_{\text{in}}}
         \;=\; -\frac{R_{f}}{R_{i}}
  \label{eq:gain}
\end{equation}
where $R_f$ is the feedback resistor and $R_i$ the input resistor.  With
$R_f=\qty{100}{\kohm}$ and $R_i=\qty{10}{\kohm}$ the theoretical gain is
$A_v = -10.0$.

\begin{figure}[htbp]
  \centering
  \begin{circuitikz}[american,scale=0.95]
    \draw (0,1.8) node[op amp] (opamp) {};
    \draw (opamp.-) to[short] ++(-1.4,0) coordinate (sum)
          to[R,l_=$R_i$\;\qty{10}{\kohm}] ++(-2.6,0)
          to[short,-o] ++(-0.4,0) node[left]{$V_{\text{in}}$};
    \draw (sum) to[short,*-] ++(0,1.1) coordinate (fb)
          to[R,l=$R_f$\;\qty{100}{\kohm}]
          (fb -| opamp.out) to[short,-*] (opamp.out);
    \draw (opamp.+) to[short] ++(0,-0.9) node[ground]{};
    \draw (opamp.out) to[short,-o] ++(1,0) node[right]{$V_{\text{out}}$};
    \node at (2.2,-2.8) {};
  \end{circuitikz}
  \caption{Inverting amplifier.\;\;
    $R_i=\qty{10}{\kohm},\;R_f=\qty{100}{\kohm}\;\Rightarrow\;A_v=-10$.}
  \label{fig:invamp}
\end{figure}

\subsection{Active Low-Pass Filter}

A unity-gain Sallen--Key low-pass filter (Figure~\ref{fig:lpf}) yields a
second-order Butterworth response when $R_1=R_2=R$ and $C_1=2C_2$.  The
transfer function is
\begin{equation}
  H(s)=\frac{1}{s^{2}R^{2}C_{1}C_{2}+2sRC_{2}+1}
  \label{eq:tf}
\end{equation}
with a cut-off frequency
$f_c = \bigl(2\pi R\sqrt{C_{1}C_{2}}\,\bigr)^{-1}$.  Using
$R=\qty{10}{\kohm}$, $C_1=\qty{22}{\nF}$, and $C_2=\qty{10}{\nF}$
gives $f_c \approx \qty{1.07}{\kHz}$.

\begin{figure}[htbp]
  \centering
  \begin{circuitikz}[american,scale=0.95]
    \draw (0,0) node[op amp] (opamp) {};
    \draw (opamp.-) to[short] ++(0,1) to[short]
          (opamp.out |- {(0,1)}) to[short] (opamp.out);
    \draw (opamp.+) to[short] ++(-1.4,0) coordinate (mid)
          to[C,l_=$C_2$\;\qty{10}{\nF}] ++(0,-1.3) node[ground]{};
    \draw (mid) to[R,l_=$R_2$\;\qty{10}{\kohm}] ++(-2.4,0) coordinate (inj)
          to[R,l_=$R_1$\;\qty{10}{\kohm},-o] ++(-2,0)
          node[left]{$V_{\text{in}}$};
    \draw (inj) to[C,l=$C_1$\;\qty{22}{\nF}]
          (inj |- opamp.out) to[short,-*] (opamp.out);
    \draw (opamp.out) to[short,-o] ++(1,0) node[right]{$V_{\text{out}}$};
    \node at (2.2,-2.6) {};
  \end{circuitikz}
  \caption{Unity-gain Sallen--Key low-pass filter.\;\;
    $f_c\approx\qty{1.07}{\kHz}$.}
  \label{fig:lpf}
\end{figure}

% ═══════════════════════════════════════════════
%  EXPERIMENTAL  SETUP
% ═══════════════════════════════════════════════
